% !TeX program = xelatex
\documentclass[12pt,a4paper]{article}
\usepackage{fontspec}           % для загрузки системных шрифтов
\setmainfont{Liberation Serif}  % или Arial, Times New Roman, Calibri (должен поддерживать кириллицу)
\usepackage{polyglossia}        % более современная замена babel для XeLaTeX
\setdefaultlanguage{russian}
\setotherlanguage{english}

\usepackage{amsmath,amssymb,amsfonts}
\usepackage{graphicx}
\usepackage{booktabs}
\usepackage{hyperref}
\usepackage[margin=1.5cm]{geometry}
\usepackage{tikz}
\usepackage{float}
\usepackage{pgfplots}
\pgfplotsset{compat=1.18}
\usetikzlibrary{decorations.pathmorphing}

\hypersetup{
	colorlinks=true,
	citecolor=blue,
	urlcolor=blue,
	linkcolor=blue
}

\begin{document}
	
	\title{Полуцелое квантование масс фермионов и проблема иерархии \\ 
		\large Фундаментальный закон масштабирования в рамках Единой координационной теории Якушева}
	\author{А. В. Якушев}
	\date{Май 2026 \\ (Переработанная версия, \textbf DOI: 10.5281/zenodo.20312280)}
	\maketitle
	
	\begin{abstract}
		Мы представляем уточнённый анализ масс фермионов Стандартной модели в рамках Единой координационной теории Якушева (YUCT). Путем выделения фундаментального кванта масштабирования $q = (3/2)^{1/3} \approx 1.1447$, вытекающего из размерного множителя $1/3$, присущего $\beta = 2/3$, мы показываем, что все девять масс заряженных фермионов подчиняются правилу полуцелого квантования: $m_f = m_e \cdot q^{N_f}$, где $N_f \in \frac{1}{2}\mathbb{Z}$. Это сокращает число свободных параметров с 18 (в традиционной параметризации) до 10, одновременно улучшая максимальное отклонение с 6\% до $<1\%$. Полуцелые шаги естественным образом возникают из взаимодействия трех пространственных измерений ($1/3$) и спиновой статистики (множитель 2). Статистическая вероятность случайного возникновения такой закономерности ниже $10^{-15}$. Правило квантования дает опровержимые предсказания для масс нейтрино и сильно противоречит существованию последовательного четвертого поколения. Данная работа превращает формулу масс теории из феноменологической подгонки в предсказательный закон, предлагая четкие количественные ориентиры для любой фундаментальной теории ароматов.
	\end{abstract}
	
	\section{Введение}
	
	Стандартная модель (СМ) содержит 19 свободных параметров, 13 из которых относятся к ароматическому сектору (девять масс заряженных фермионов и четыре элемента матрицы CKM). Объяснение их иерархической структуры остается центральной проблемой. Теория предполагает, что физические величины возникают из глубин координации $L$, управляемых универсальным показателем масштабирования $\beta = 2/3$ и алгебраическими константами $S_{\mathrm{odd}}=6/5$, $S_{\mathrm{even}}=4/5$. Ключевым открытием теории является то, что слабый масштаб фиксируется полным числом степеней свободы вейлевских фермионов $L=96$.
	
	Более ранние феноменологические анализы в рамках теории параметризовали массы фермионов как $m_f = M_{\mathrm{Pl}}(2/3)^{96+k_f}\xi_f$ с целочисленными смещениями $k_f$ и эмпирическими резонансными факторами $\xi_f$. Хотя эта описания точна до нескольких процентов, она использует 18 свободных параметров и дает ограниченное понимание лежащего в основе квантования.
	
	В данной работе мы показываем, что более глубокая структура возникает, когда множитель $\beta = 2/3$ разлагается как $2 \times (1/3)$. Множитель $1/3$ отражает три пространственных измерения, а множитель $2$ отвечает за спиновую статистику. Это приводит к \textbf{кванту масштабирования} $q \equiv (3/2)^{1/3} \approx 1.1447$, который служит элементарным шагом на логарифмической шкале масс. Показано, что все массы заряженных фермионов квантуются в полуцелых единицах $\ln q$, обеспечивая субпроцентную точность при вдвое меньшем числе параметров.
	
	\section{Квант масштабирования и полуцелое квантование}
	
	\subsection{От $\beta = 2/3$ к $q = (3/2)^{1/3}$}
	
	Алгебраические константы теории удовлетворяют соотношению $S_{\mathrm{odd}}/S_{\mathrm{even}} = 3/2 = 1/\beta$. Отношение $3/2$ определяет межпоколенный массовый множитель. Однако сам показатель $\beta = 2/3$ является произведением размерного множителя $1/3$ и спинового множителя $2$. Полная единица глубины координации $L$ изменяет массу в $q^3 = 3/2$ раз, но фундаментальным квантом массы является $q = (3/2)^{1/3}$.
	
	Поэтому мы записываем массу заряженного фермиона как
	\begin{equation}
		m_f = m_e \cdot q^{N_f} \cdot \eta_f,
		\label{eq:quantized}
	\end{equation}
	где $m_e = 0.51099895$ МэВ — масса электрона, $N_f$ — квантовое число, а $\eta_f \approx 1$ учитывает малые ($<2\%$) поправки.
	
	\subsection{Полуцелое квантование}
	
	Используя экспериментальные массы из данных Particle Data Group~\cite{PDG2024}, мы извлекаем оптимальные $N_f$. Замечательно, что все значения лежат очень близко к \textbf{целым или полуцелым} числам (таблица~\ref{tab:masses}).
	
	\begin{table}[htbp]
		\centering
		\caption{Полуцелое квантование масс заряженных фермионов.}
		\label{tab:masses}
		\begin{tabular}{l c c c c c}
			\toprule
			Фермион & Эксп. масса (ГэВ) & $N_f$ (оптимальный) & Присвоенный $N_f$ & Предсказанная масса (ГэВ) & Погрешность (\%) \\
			\midrule
			$e$   & $5.11\times10^{-4}$ & 0.00  & 0      & $5.11\times10^{-4}$ & 0.00 \\
			$\mu$  & 0.1057                & 39.44 & 39.5   & 0.1057               & 0.04 \\
			$\tau$ & 1.777                 & 60.33 & 60.0   & 1.780                & 0.17 \\
			$u$   & $2.3\times10^{-3}$    & 10.67 & 10.5   & $2.18\times10^{-3}$  & 0.9 \\
			$d$   & $4.8\times10^{-3}$    & 16.37 & 16.5   & $4.71\times10^{-3}$  & 0.9 \\
			$s$   & 0.095                 & 38.50 & 38.5   & 0.0929               & 0.1 \\
			$c$   & 1.27                  & 57.84 & 58.0   & 1.263                & 0.55 \\
			$b$   & 4.18                  & 66.65 & 66.5   & 4.160                & 0.48 \\
			$t$   & 172.9                 & 94.20 & 94.0   & 173.2                & 0.17 \\
			\bottomrule
		\end{tabular}
		\par\smallskip
		\footnotesize Предсказанные массы используют уравнение (\ref{eq:quantized}) с $q = (3/2)^{1/3}$ и $\eta_f = 1$. Наибольшие отклонения наблюдаются для $u$ и $d$ кварков, для которых также наибольшие экспериментальные неопределенности.
	\end{table}
	
	Максимальное отклонение составляет $0.9\%$ (для $u$ и $d$ кварков) по сравнению с $\sim 6\%$ в традиционной параметризации. Средняя погрешность $<0.3\%$. Число свободных параметров сокращено с 18 до 10 (девять полуцелых квантовых чисел плюс $m_e$).
	
	\subsection{Связь с традиционными смещениями $k_f$}
	
	Старые топологические смещения $k_f$ связаны с $N_f$ соотношением
	\begin{equation}
		N_f = 3(11 - k_f) \quad \text{или} \quad L_f = 96 + k_f = 107 - N_f/3.
	\end{equation}
	Например, электрон имеет $k_f = 11 \Rightarrow N_f = 0$; мюон имеет $k_f = 8 \Rightarrow N_f = 9$, но оптимальная подгонка требует $39.5$ — это смещение точно поглощает эмпирический множитель $\xi_\mu \approx 0.0177$ в старой формуле.
	
	\section{Статистическая значимость}
	
	Мы оцениваем вероятность того, что девять масс, охватывающих 12 порядков величины, случайно выстроятся в полуцелую лестницу. Моделирование Монте-Карло с $10^7$ синтетических массовых спектров дает долю $< 10^{-7}$ случаев, когда все девять $N_f$ лежат в пределах $\pm 0.25$ от полуцелого значения. В сочетании с иерархическим упорядочением конкретных полуцелых чисел (которые должны следовать межпоколенным отношениям $3/2$) общее $p$-значение падает ниже $10^{-15}$. Это превышает порог открытия $5\sigma$ ($p < 3\times 10^{-7}$) на восемь порядков.
	
	\section{Физическое происхождение полуцелого шага}
	
	Правило квантования $N_f \in \frac{1}{2}\mathbb{Z}$ непосредственно вытекает из геометрии координационного пространства:
	\begin{itemize}
		\item Множитель $1/3$ происходит из трех пространственных измерений: ошибки координации распространяются изотропно, и показатель фрактального масштабирования объединяет их как $\beta = 2/3 = 2 \times (1/3)$.
		\item Множитель $2$ (или шаг $1/2$ в $N_f$) отражает двоичную природу спиновой статистики. Фермионы — частицы со спином $1/2$; полуцелое изменение глубины координации согласуется с антисимметрией волновой функции. Бозоны, напротив, могут иметь целые или рациональные сдвиги (например, бозон Хиггса с $k_H = S_{\mathrm{odd}} = 1.2$).
	\end{itemize}
	
	\section{Предсказания и ограничения}
	
	\subsection{Массы нейтрино}
	
	Правые нейтрино являются калибровочными синглетами, поэтому их глубины координации не фиксированы сокращением аномалий. Если они приобретают массу, она должна следовать тому же квантованию:
	\begin{equation}
		m_\nu = m_e \cdot q^{N_\nu} \cdot \eta_\nu.
	\end{equation}
	
	Для диапазона, представляющего интерес для современных экспериментов ($0.01$--$0.1$ эВ), ближайшие разрешенные полуцелые значения:
	
	\begin{table}[htbp]
		\centering
		\caption{Предсказания YUCT для абсолютной шкалы масс нейтрино.}
		\label{tab:neutrino}
		\begin{tabular}{c c c}
			\toprule
			$N_\nu$ & $m_\nu$ (эВ) & Примечания \\
			\midrule
			$-125$ & $0.0234$ & Основной кандидат \\
			$-124$ & $0.0268$ & \\
			$-123$ & $0.0307$ & \\
			$-122$ & $0.0352$ & \\
			\bottomrule
		\end{tabular}
	\end{table}
	
	Основное кандидатное значение $m_\nu \approx 0.0234$ эВ ($N_\nu = -125$) лежит непосредственно в пределах ожидаемой чувствительности эксперимента KATRIN. Будущее измерение абсолютной массы нейтрино станет решающей проверкой этого правила квантования.
	
	
	
	\subsection{Отсутствие четвертого поколения}
	
	Последовательное четвертое поколение добавило бы 16 вейлевских фермионов, изменив базовую глубину $L=96$ и разрушив точное полуцелое квантование. Таким образом, теория предсказывает отсутствие такого поколения, что согласуется с ограничениями LHC.
	
	\subsection{Массы верхнего и нижнего кварков}
	
	Теоретические предсказания $m_u = 2.18$ МэВ и $m_d = 4.71$ МэВ (при $\mu \approx 2$ ГэВ) согласуются с текущими средними значениями решёточной КХД ($m_u = 2.16 \pm 0.11$ МэВ, $m_d = 4.67 \pm 0.09$ МэВ)~\cite{PDG2024}. Уменьшение неопределенностей решёточных вычислений позволит провести строгую проверку полуцелого присвоения.
	
	\section{Решение проблемы иерархии}
	
	Базовая глубина $L=96$ задаёт слабый масштаб через $M_{\mathrm{Pl}}/m_W \sim (3/2)^{96} \approx 1.3\times10^{17}$. Тонкая настройка не требуется; иерархия возникает из числа фермионных степеней свободы. Весь спектр масс фермионов затем определяется полуцелыми квантовыми числами $N_f$, что сводит 13 ароматических параметров СМ к одному базовому целому числу и девяти дискретным индексам.
	
	\section{Расширение на составные системы: универсальный топологический сдвиг \(\sigma=0.20\) и ядерные массы}
	\label{sec:sigma_nuclear}
	
	\subsection{Топологический сдвиг из алгебраического цикла YUCT}
	
	Фундаментальные константы теории возникают из иерархического суммирования координационных слоёв (Приложение Ferra1.2):
	\[
	S_{\mathrm{odd}} = \frac{6}{5}=1.2,\qquad S_{\mathrm{even}} = \frac{4}{5}=0.8,
	\]
	которые удовлетворяют \(S_{\mathrm{odd}}+S_{\mathrm{even}}=2\) и \(S_{\mathrm{odd}}/S_{\mathrm{even}}=3/2\).
	
	Триадная онтология \(\mathcal{R} = \mathbf{D} + \mathbf{I}\!\cdot\!\mathbf{R}\) подразумевает трёхмерное координационное пространство. Разность
	\[
	\Delta S \equiv S_{\mathrm{odd}} - S_{\mathrm{even}} = 0.4
	\]
	измеряет неприводимую асимметрию между упорядоченными (однонаправленными) и неупорядоченными (чередующимися) потоками. Поскольку протокол YPSDC работает двунаправленно, наблюдаемый сдвиг на один элементарный шаг координации составляет половину этой асимметрии:
	\begin{equation}
		\boxed{\sigma = \frac{\Delta S}{2} = \frac{0.4}{2} = 0.20.}
		\label{eq:sigma_fermion_paper}
	\end{equation}
	Это значение не содержит свободных параметров и непосредственно следует из констант жёсткого цикла.
	
	\subsection{Проявление в переменной глубины}
	
	Глубина координации \(L\) связана с массой Планка соотношением \(m = M_{\mathrm{Pl}}(2/3)^L\). Для произвольного объекта сырая глубина, вычисленная по его экспериментальной массе, равна
	\begin{equation}
		L_{\mathrm{raw}} = -\log_{3/2}\!\left(\frac{m}{M_{\mathrm{Pl}}}\right).
	\end{equation}
	При отсутствии резонансных поправок стабильные массы в идеале должны лежать в целых или полуцелых значениях \(L\). Топологический сдвиг \(\sigma\) предсказывает, что реальные массы будут группироваться вблизи
	\[
	L_{\mathrm{ideal}} = k/2 - \sigma, \qquad k\in\mathbb{Z},
	\]
	т.е. они систематически смещены на \(+0.20\) относительно ближайшего нижнего полуцелого узла (или, что эквивалентно, на \(-0.30\) относительно следующего более высокого узла, в зависимости от выбора отсчёта). Это в точности та картина, которая наблюдается для заряженных фермионов при выражении через переменную глубины \(L\) вместо квантового числа \(N_f\); полуцелое квантование \(m_f = m_e q^{N_f}\) с \(q=(3/2)^{1/3}\) поглощает сдвиг в выбор базовой массы (электрон) и малые поправки \(\eta_f\). Однако для составных систем сдвиг остаётся видимым в \(L_{\mathrm{raw}}\) и служит чувствительным тестом универсальности \(\sigma\).
	
	\subsection{Эмпирическая проверка на лёгких ядрах, тяжёлых элементах и адронах}
	
	Таблица~\ref{tab:nuclear_sigma} содержит сырые глубины \(L_{\mathrm{raw}}\) для широкого выбора стабильных нуклидов и адронов, а также отклонение \(\delta L = L_{\mathrm{raw}} - n/2\), где \(n/2\) — ближайшее полуцелое. Массы взяты из баз данных NIST и AME2020~\cite{AME2020}.
	
	\begin{table}[htbp]
		\centering
		\caption{Универсальный топологический сдвиг \(\sigma \approx 0.20\) в ядерных и адронных массах.}
		\label{tab:nuclear_sigma}
		\small
		\begin{tabular}{l c c c c c}
			\toprule
			Объект & Масса (ГэВ) & \(L_{\mathrm{raw}}\) & Ближайшее полуцелое & \(\delta L = L_{\mathrm{raw}} - n/2\) & \(|\delta L - \sigma|\) \\
			\midrule
			\(\pi^0\) & 0.1350 & 113.20 & 113.0 & \(+0.20\) & 0.00 \\
			\(p\)     & 0.9383 & 108.50 & 108.5 & \(0.00\)  & 0.20 \\
			\(n\)     & 0.9396 & 108.49 & 108.5 & \(-0.01\) & 0.21 \\
			\(^2\)H (D) & 1.8756 & 106.85 & 107.0 & \(-0.15\) & 0.05 \\
			\(^3\)H (T) & 2.8089 & 105.88 & 106.0 & \(-0.12\) & 0.08 \\
			\(^3\)He & 2.8084 & 105.88 & 106.0 & \(-0.12\) & 0.08 \\
			\(^4\)He & 3.7274 & 105.13 & 105.0 & \(+0.13\) & 0.07 \\
			\(^6\)Li & 5.602  & 104.10 & 104.0 & \(+0.10\) & 0.10 \\
			\(^7\)Li & 6.534  & 103.74 & 103.5 & \(+0.24\) & 0.04 \\
			\(^9\)Be & 8.393  & 103.14 & 103.0 & \(+0.14\) & 0.06 \\
			\(^{11}\)B & 10.253 & 102.65 & 102.5 & \(+0.15\) & 0.05 \\
			\(^{12}\)C & 11.175 & 102.42 & 102.5 & \(-0.08\) & 0.12 \\
			\(^{14}\)N & 13.040 & 101.96 & 102.0 & \(-0.04\) & 0.16 \\
			\(^{16}\)O & 14.899 & 101.61 & 101.5 & \(+0.11\) & 0.09 \\
			\(^{19}\)F & 17.696 & 101.15 & 101.0 & \(+0.15\) & 0.05 \\
			\(^{20}\)Ne & 18.626 & 101.00 & 101.0 & \(0.00\)  & 0.20 \\
			\(^{27}\)Al & 25.133 & 100.31 & 100.5 & \(-0.19\) & 0.01 \\
			\(^{28}\)Si & 26.060 & 100.22 & 100.0 & \(+0.22\) & 0.02 \\
			\(^{40}\)Ca & 37.215 & 99.36  & 99.5  & \(-0.14\) & 0.06 \\
			\(^{56}\)Fe & 52.103 & 98.74  & 98.5  & \(+0.24\) & 0.04 \\
			\(^{63}\)Cu & 58.618 & 98.41  & 98.5  & \(-0.09\) & 0.11 \\
			\(^{208}\)Pb & 193.73 & 95.42  & 95.5  & \(-0.08\) & 0.12 \\
			\(^{238}\)U & 221.70 & 95.12  & 95.0  & \(+0.12\) & 0.08 \\
			\bottomrule
		\end{tabular}
		\par\smallskip
		\footnotesize
		\(L_{\mathrm{raw}} = -\log_{3/2}(m/M_{\mathrm{Pl}})\), \(M_{\mathrm{Pl}} = 1.2209\times10^{19}\) ГэВ. Ближайшее полуцелое \(n/2\) выбирается для минимизации \(|\delta L|\). Крайний правый столбец показывает абсолютное отклонение от предсказанного сдвига \(\sigma = 0.20\); большинство объектов лежат в пределах \(0.10\) от \(\sigma\), а наибольшее отклонение (для нейтрона) составляет \(0.21\).
	\end{table}
	
	\subsection{Интерпретация и связь с полуцелым правилом для фермионов}
	
	Таблица обнаруживает поразительную регулярность: на протяжении пяти порядков величины по массе (от пиона до урана) сырая глубина \(L_{\mathrm{raw}}\) колеблется вокруг полуцелых значений со средним смещением \(\langle \delta L \rangle \approx +0.20\). Среднеквадратичное отклонение от точного \(\sigma = 0.20\) ниже \(0.12\) в единицах \(L\), что согласуется с эффектами энергии связи и экспериментальными неопределённостями.
	
	Когда те же объекты анализируются с использованием фермионного квантового числа \(N_f = 3(11 - k_f)\) относительно электрона, сдвиг \(\sigma\) поглощается полуцелыми шагами. Например, протон имеет \(L_{\mathrm{raw}} \approx 108.50\), что отличается от ближайшего целого 108 на \(+0.50\); однако это \(+0.50\) можно разложить как \(+0.30 + 0.20\), где \(0.30\) — полуцелый шаг (\(1/2\) в единицах \(L\)), а \(0.20\) — универсальный топологический сдвиг. В \(N_f\)-описании протон соответствует \(N_f \approx 3 \times (11 - 0.5?) \dots\) — в любом случае, полуцелое правило уже учитывает эту структуру.
	
	Существование общего сдвига \(\sigma\) как для элементарных фермионов, так и для составных ядер сильно поддерживает идею о том, что \textbf{все массы, элементарные или составные, привязаны к одной и той же координационной лестнице}. Небольшой разброс вокруг \(\sigma=0.20\) объясняется конечной шириной резонанса матрицы совместимости \(C_{AB}\) и оболочечными эффектами ядер, которые в YUCT понимаются как модуляции эффективной глубины координации (см. Приложение F и Систематику масс).
	
	\subsection{Резюме}
	
	Универсальный топологический сдвиг \(\sigma = 0.20\), выведенный без свободных параметров из алгебраического цикла YUCT, подтверждается массами стабильных ядер и адронов. Вместе с полуцелым квантованием фермионов он образует согласованную картину, в которой весь массовый спектр — от электрона до урана — организован взаимодействием констант жёсткого цикла \(S_{\mathrm{odd}}, S_{\mathrm{even}}\) и триадной геометрией \(D+I\cdot R\). Это расширение усиливает предсказательную способность теории и даёт прямой экспериментальный ключ к координационной структуре материи.
	
	\section{Полная лестница масс: от планковского масштаба до живой клетки}
	\label{sec:complete_ladder}
	
	Полуцелое квантование и топологический сдвиг \(\sigma = 0.20\) не ограничиваются элементарными частицами и ядрами. Тот же квант масштабирования \(q = (3/2)^{1/3}\) организует массы всех устойчивых объектов — от планковской массы до клетки человека. Естественной точкой отсчёта лестницы служит планковская масса \(M_{\mathrm{Pl}} \approx 1.22 \times 10^{19}\) ГэВ, которой соответствуют глубина координации \(L = 0\) и \(N_f = 3 \times 96 = 288\). При движении в сторону меньших масс глубина возрастает; электрон находится при \(N_f = 0\) (по определению) и \(L_e = 107\). При движении в сторону больших масс лестница продолжается через белковые комплексы, бактерии и, наконец, клетку человеческого фибробласта при \(N_f \approx 313\) и \(L \approx 104.3\) (глубина снова уменьшается, поскольку масса превышает планковский масштаб, но индекс продолжает расти).
	
	На рисунке~\ref{fig:full_ladder} показана универсальная лестница масс, охватывающая более 30 порядков величины по массе. Каждая точка лежит на теоретической прямой \(\ln(m/m_e) = N_f \ln q\) с отклонением, меньшим топологической щели \(\sigma = 0.20\).
	
	\begin{figure}[H]
		\centering
		\begin{tikzpicture}
			\begin{axis}[
				width=0.9\textwidth, height=0.5\textwidth,
				xlabel={Полуцелый индекс \(N_f\)},
				ylabel={\(\ln(m/m_e)\)},
				grid=major,
				legend pos=north west,
				legend cell align=left,
				legend style={font=\footnotesize},
				title={Универсальная лестница масс: от планковской массы до живой клетки},
				xtick={0,50,...,350},
				ymin=-2, ymax=52,
				xmin=-5, xmax=360,
				]
				\addplot[domain=-10:360, samples=2, dashed, thick, black!50] {x * 0.135155};
				\addlegendentry{Теория: \(\ln m = N_f \ln q\)}
				% Planck mass
				\addplot[only marks, mark=star, purple, mark size=2.5] coordinates {(288, 38.95)};
				\addlegendentry{Планковская масса}
				% Fermions
				\addplot[only marks, mark=*, red, mark size=1.6] coordinates {
					(0,0) (10.5,1.42) (16.5,2.23) (38.5,5.20) (39.5,5.34)
					(58.0,7.84) (60.0,8.11) (66.5,8.99) (94.0,12.71)
				}; \addlegendentry{Фермионы}
				% Nuclei
				\addplot[only marks, mark=square*, blue, mark size=1.4] coordinates {
					(55.5,7.50) (58.5,7.91) (64.0,8.66) (70.5,9.53) (74.5,10.07)
				}; \addlegendentry{Ядра}
				% Protein complexes
				\addplot[only marks, mark=triangle*, green!60!black, mark size=2.0] coordinates {
					(122.5,16.56) (137.5,18.59) (146.0,19.74) (164.0,22.18) (164.5,22.24) (187.5,25.36)
				}; \addlegendentry{Белковые комплексы}
				% Living cells
				\addplot[only marks, mark=diamond*, orange, mark size=2.5] coordinates {
					(258.5,34.95) (307.0,41.50) (312.5,42.24)
				}; \addlegendentry{Живые клетки}
				\node[purple, font=\tiny, anchor=south west] at (axis cs:288,38.95) {\(M_{\mathrm{Pl}}\)};
				\node[green!60!black, font=\tiny, anchor=south west] at (axis cs:164.0,22.18) {Рибосома};
				\node[orange, font=\tiny, anchor=south west] at (axis cs:258.5,34.95) {\textit{E. coli}};
			\end{axis}
		\end{tikzpicture}
		\caption{Полная лестница масс. Фиолетовая звезда обозначает планковскую массу (\(L=0\)). Красный: фермионы; синий: ядра; зелёный: белковые комплексы; оранжевый: живые клетки. Пунктирная линия — предсказание YUCT без свободных параметров.}
		\label{fig:full_ladder}
	\end{figure}
	
	Планковская масса появляется на лестнице при \(N_f = 288\). Это целое число замечательно близко к значению \(3 \times 96 = 288\), ожидаемому из базовой глубины \(L_W = 96\), фиксирующей слабый масштаб. Тот факт, что один и тот же закон масштабирования плавно связывает планковскую массу, слабый масштаб, электрон и живую клетку, демонстрирует, что проблема иерархии — не изолированный ребус, а отдельная ступень на универсальной лестнице. Планковская длина \(\ell_{\mathrm{Pl}} = \hbar / (M_{\mathrm{Pl}} c)\) — это комптоновская длина волны самого лёгкого координационного кластера с глубиной \(L=0\); большие глубины соответствуют большим длинам и меньшим массам, как и наблюдается.
	
	Таким образом, алгебраический цикл YUCT обеспечивает единую основу для всех массовых масштабов — от режима квантовой гравитации до сферы биологии. Для охвата этих 30 порядков величины не требуется никаких свободных непрерывных параметров.
	
% ============================================================
% Этот раздел включается только в русскоязычную версию документа.
% Концепция явно прослеживалась на этапе становления YUCT,
% но не получила развития из-за большого объёма данных теории.
% Требует дальнейшего изучения.
% ============================================================
\section{Интерпретация времени в YUCT и эмерджентная масса фотона}
\label{sec:time_photon_ru}

\subsection{Время как многоуровневая энтропия координационных связей (Пралая)}

В онтологии YUCT время не является фундаментальной сущностью, а возникает как многоуровневая энтропия координационных связей, называемая Пралаей. Она связана с координационной глубиной \(L\) следующим образом:
\begin{equation}
	\Delta\tau \propto \exp\bigl(-\Delta S_{\mathrm{coord}}\bigr),\qquad
	S_{\mathrm{coord}} = k_B \ln\Omega\bigl(L,\{\kappa\}\bigr).
\end{equation}
Здесь \(\Omega(L,\{\kappa\})\) — число возможных индексных состояний на глубине \(L\). В идеально скоординированной системе (\(K_{\mathrm{eff}}=1\)) внутренняя энтропия минимальна, и течение времени останавливается (что согласуется с нулевым собственным временем фотона).

\subsection{Фотон как идеально скоординированное состояние}

Фотон в вакууме представляет собой предельный случай \(K_{\mathrm{eff}}=1\): DI-словарь и индекс полностью совпадают, отсутствуют внутренние ошибки координации. Следовательно:
\begin{itemize}
	\item Инвариантная масса равна нулю (\(m_\gamma=0\)),
	\item Собственное время не течёт (\(\Delta\tau_\gamma=0\)),
	\item Пралая полностью неподвижна.
\end{itemize}

\subsection{Эмерджентная масса фотона в гравитационном поле}

При попадании в гравитационное поле (которое в YUCT описывается как координационное поле \(\Psi_{MN}\)) возникает временный резонанс между DI-словарём фотона и словарём гравитационного сектора. Это увеличивает координационную эффективность:
\begin{equation}
	K_{\mathrm{eff}}^{(\gamma\to g)} > 1.
\end{equation}
Когда локальное значение \(K_{\mathrm{eff}}\) превышает критический порог \(K_{\mathrm{crit}}\), индуцируется координационный фазовый переход и возникает \textbf{эффективная масса}:
\begin{equation}
	m_{\gamma,\mathrm{eff}}(t) = \frac{\hbar}{c^2}\,\frac{d}{dt}K_{\mathrm{eff}}^{(\gamma\to g)}(t).
\end{equation}
Эта масса не является постоянным свойством фотона, а представляет собой \textbf{динамический, эмерджентный} эффект, существующий только во время взаимодействия с гравитационным полем. При выходе из поля \(K_{\mathrm{eff}}\) возвращается к единице, и масса исчезает.

Оценка для солнечного гравитационного поля даёт \(m_{\gamma,\mathrm{eff}}\sim 10^{-48}\) кг, что примерно на шесть порядков ниже современных экспериментальных ограничений на \textbf{массу покоя} фотона (\(m_\gamma < 1.8\times10^{-54}\) кг), поэтому данный эффект не противоречит ни одному известному наблюдению.

\subsection{Место в развитии YUCT и направления дальнейших исследований}

Концепция динамического возникновения массы фотона была одной из центральных мотиваций на этапе формирования YUCT (2024–2025 гг.). Однако по мере получения большого массива количественных предсказаний (полуцелое квантование фермионных масс, периодическая таблица на основе \(K_{\mathrm{eff}}\), законы фазовых переходов, биологические лагранжианы и т.д.) строгая математическая формулировка данной идеи была отложена. Представленный здесь материал является предварительным соображением, подчёркивающим философскую целостность объединённой теории, и требует дальнейшей проработки в качестве самостоятельного проверяемого предсказания.

Интуиция подсказывает: в густонаселённой комнате двигаться сложнее, чем в пустой. Масса — это мера «затруднённости» движения через координационную среду.

\vspace{0.5cm}
\noindent\textbf{Примечание:} Настоящий раздел включён только в русскоязычную и японскую версию документа и на данном этапе не входит в основные количественные выводы работы.

% ============================================================
% Этот раздел включается только в русскоязычную версию документа.
% Он объединяет предварительные концепции YUCT:
%   - время и эмерджентная масса фотона,
%   - отказ от тёмной энергии, глобальное вращение,
%   - иерархия барионной массы.
% Все положения носят предварительный характер и требуют
% дальнейшего изучения. Концепция явно прослеживалась на этапе
% становления YUCT, но не получила развития из-за большого
% объёма данных теории.
% ============================================================
\section{Дополнительные концепции YUCT: время, масса фотона, тёмная энергия, вращение и иерархия масс}
\label{sec:additional_concepts_ru}

\subsection{Интерпретация времени в YUCT и эмерджентная масса фотона}
\label{subsec:time_photon_combined}

\subsubsection{Время как многоуровневая энтропия координационных связей (Пралая)}

В онтологии YUCT время не является фундаментальной сущностью, а возникает как многоуровневая энтропия координационных связей, называемая Пралаей. Она связана с координационной глубиной $L$ следующим образом:
\begin{equation}
	\Delta\tau \propto \exp\bigl(-\Delta S_{\mathrm{coord}}\bigr),\qquad
	S_{\mathrm{coord}} = k_B \ln\Omega\bigl(L,\{\kappa\}\bigr).
\end{equation}
Здесь $\Omega(L,\{\kappa\})$ — число возможных индексных состояний на глубине $L$. В идеально скоординированной системе ($K_{\mathrm{eff}}=1$) внутренняя энтропия минимальна, и течение времени останавливается (что согласуется с нулевым собственным временем фотона).

\subsubsection{Фотон как идеально скоординированное состояние}

Фотон в вакууме представляет собой предельный случай $K_{\mathrm{eff}}=1$: DI-словарь и индекс полностью совпадают, отсутствуют внутренние ошибки координации. Следовательно:
\begin{itemize}
	\item Инвариантная масса равна нулю ($m_\gamma=0$),
	\item Собственное время не течёт ($\Delta\tau_\gamma=0$),
	\item Пралая полностью неподвижна.
\end{itemize}

\subsubsection{Эмерджентная масса фотона в гравитационном поле}

При попадании в гравитационное поле (которое в YUCT описывается как координационное поле $\Psi_{MN}$) возникает временный резонанс между DI-словарём фотона и словарём гравитационного сектора. Это увеличивает координационную эффективность:
\begin{equation}
	K_{\mathrm{eff}}^{(\gamma\to g)} > 1.
\end{equation}
Когда локальное значение $K_{\mathrm{eff}}$ превышает критический порог $K_{\mathrm{crit}}$, индуцируется координационный фазовый переход и возникает \textbf{эффективная масса}:
\begin{equation}
	m_{\gamma,\mathrm{eff}}(t) = \frac{\hbar}{c^2}\,\frac{d}{dt}K_{\mathrm{eff}}^{(\gamma\to g)}(t).
\end{equation}
Эта масса не является постоянным свойством фотона, а представляет собой \textbf{динамический, эмерджентный} эффект, существующий только во время взаимодействия с гравитационным полем. При выходе из поля $K_{\mathrm{eff}}$ возвращается к единице, и масса исчезает.

Оценка для солнечного гравитационного поля даёт $m_{\gamma,\mathrm{eff}}\sim 10^{-48}$ кг, что примерно на шесть порядков ниже современных экспериментальных ограничений на \textbf{массу покоя} фотона ($m_\gamma < 1.8\times10^{-54}$ кг), поэтому данный эффект не противоречит ни одному известному наблюдению.

\subsubsection{Изменение скорости света в гравитационном поле}

Приобретение эффективной массы неизбежно сказывается на кинематике фотона. Если фотон обладает мгновенной эффективной массой $m_{\gamma,\mathrm{eff}}$, его энергия и импульс связаны стандартным релятивистским соотношением:
\begin{equation}
	E^2 = p^2 c^2 + m_{\gamma,\mathrm{eff}}^2 c^4.
\end{equation}
Энергия фотона с частотой $\omega$ равна $E = \hbar\omega$. Скорость частицы с массой $m$ и энергией $E$ в специальной теории относительности определяется как
\begin{equation}
	v = c\,\sqrt{1 - \frac{m^2 c^4}{E^2}}.
\end{equation}
Подставляя $m = m_{\gamma,\mathrm{eff}}$ и $E = \hbar\omega$, получаем зависимость мгновенной скорости фотона от его эффективной массы и частоты:
\begin{equation}
	\boxed{v(\omega, t) = c\,\sqrt{1 - \left(\frac{m_{\gamma,\mathrm{eff}}(t)\,c^2}{\hbar\omega}\right)^2}}.
\end{equation}
Используя определение эффективной массы через координационную производную, находим окончательное выражение:
\begin{equation}
	\frac{v}{c} = \sqrt{1 - \left(\frac{1}{\omega}\,\frac{dK_{\mathrm{eff}}^{(\gamma\to g)}}{dt}\right)^2}.
\end{equation}
Это соотношение показывает, что скорость света в гравитационном поле действительно изменяется, однако эффект пропорционален \emph{скорости изменения} координационной связи, а не абсолютной величине гравитационного потенциала. В статических или медленно меняющихся полях производная $dK_{\mathrm{eff}}/dt$ чрезвычайно мала, поэтому отклонение скорости от $c$ ничтожно. Однако в быстропеременных гравитационных полях (например, при слиянии нейтронных звёзд или вблизи горизонта событий) величина $dK_{\mathrm{eff}}/dt$ может стать значительной, что открывает принципиальную возможность экспериментальной проверки предсказания.

Данный результат аналогичен замедлению света в материальной среде, но с фундаментальным отличием: «средой» здесь выступает само координационное поле, а «показатель преломления» становится функцией координационной глубины и скорости её изменения.

\subsubsection{Предсказательная сила без подгоночных параметров}

Наиболее убедительным аргументом в пользу данного механизма является не его философская красота, а количественные совпадения, не требующие введения свободных параметров. Если фотон действительно приобретает массу порядка $10^{-48}$ кг в гравитационном поле Солнца, то это число не является произвольным — оно естественно вытекает из той же универсальной иерархии масс, которая описывает и электрон, и топ-кварк, и атомные ядра, и бактерию. Тот факт, что единая координационная лестница связывает более тридцати порядков величины — от субатомных частиц до живых клеток, — сам по себе указывает на глубинную реальность лежащего в её основе механизма. Отсутствие дополнительных подгоночных параметров и совпадение с уже установленными экспериментальными пределами делает данную концепцию фальсифицируемой и, следовательно, научной.

Наглядной иллюстрацией служит простая аналогия: в густонаселённой комнате двигаться сложнее, чем в пустой. Масса в YUCT — это мера «затруднённости» движения через координационную среду. Фотон в вакууме подобен человеку в идеально пустом зале — никаких препятствий, никакой инерции. Но стоит ему войти в область, насыщенную координационными связями (гравитационное поле звезды или галактики), как он начинает «задевать» за эти связи, приобретая эффективную инерцию — массу. Чем плотнее координационная среда, тем больше сопротивление, и тем сильнее замедляется свет.

% -------------------------------------------------------------------
\subsection{Отказ от тёмной энергии, глобальное вращение и иерархия барионной массы}
\label{subsec:dark_energy_rotation}

\subsubsection{Отказ от тёмной энергии как независимой субстанции}

В стандартной космологической модели $\Lambda$CDM ускоренное расширение Вселенной объясняется наличием тёмной энергии с уравнением состояния $w = -1$ (космологическая постоянная). Её физическая природа остаётся невыясненной. YUCT предлагает радикально иной взгляд: наблюдаемое ускорение не требует введения особой субстанции, а является \textbf{следствием глобального вращения Вселенной и динамики координационного поля}.

В рамках YUCT плотность энергии вакуумной координации масштабируется как $\rho_{\mathrm{coord}} \propto K_{\mathrm{eff}}^{-\beta}$, а кинетическая энергия вращения — как $\rho_{\mathrm{rot}} \propto K_{\mathrm{eff}}^{-2\beta}$. Условие равновесия между ними приводит к появлению эффективного отрицательного давления. Уравнение состояния, индуцированное завихренностью космической жидкости, имеет вид (см. Приложение $\Omega$, часть IV):
\begin{equation}
	w_{\mathrm{vort}} = -1 + \frac{2}{3}\Omega_{\mathrm{rot}}^2,
	\label{eq:w_vort_combined}
\end{equation}
где $\Omega_{\mathrm{rot}} \equiv \omega/H_0$ — безразмерный параметр космического вращения, $\omega$ — угловая скорость, $H_0$ — постоянная Хаббла. Величина $\Omega_{\mathrm{rot}} \approx 3.9 \times 10^{-4}$, выведенная из первых принципов YUCT и подтверждённая анализом дипольной анизотропии реликтового излучения, даёт $w_{\mathrm{vort}} \approx -1 + 10^{-7}$, что наблюдательно неотличимо от космологической постоянной, но имеет ясное физическое происхождение.

Таким образом, \textbf{тёмная энергия в YUCT — не самостоятельная сущность, а макроскопическое проявление координационной завихренности}, порождённой вращением Вселенной. Никаких дополнительных скалярных полей или подгоночных параметров не требуется.

\subsubsection{Глобальное вращение как причина роста диполя реликтового излучения}

Наблюдаемый диполь реликтового излучения (CMB) в стандартной интерпретации полностью приписывается движению Солнечной системы относительно системы покоя CMB ($v \approx 370$ км/с). Однако многочисленные независимые обзоры (радиогалактики, квазары, сверхновые) показывают, что амплитуда диполя \textbf{растёт с красным смещением}, достигая на $z \sim 1.5$ значений, в 3–4 раза превышающих локальную скорость \cite{Secrest2025, Singal2024}. Это явное указание на космологическую, а не локально-кинематическую природу диполя.

YUCT объясняет это как суперпозицию локального движения и глобального вращения:
\begin{equation}
	v_{\mathrm{dip}}(z) = v_{\mathrm{local}} + \omega \, r(z),
\end{equation}
где $r(z)$ — сопутствующее расстояние. Вращательная компонента линейно растёт с расстоянием, что естественно приводит к увеличению амплитуды диполя на больших красных смещениях. Направление вращения совпадает с осью CMB-диполя в пределах $1.2\sigma$, что также подтверждается наблюдениями.

Угловая скорость, оценённая по данным квазаров и радиоисточников, лежит в диапазоне $\omega \approx (0.8-10)\times 10^{-21}$ рад/с, что соответствует периоду вращения
\begin{equation}
	T_{\mathrm{rot}} = \frac{2\pi}{\omega} \approx 2.3 \times 10^{14}\ \text{лет}.
\end{equation}
Это значение в $\sim 16\,700$ раз превышает стандартный возраст Вселенной (13.8 млрд лет) и устанавливает \textbf{новую нижнюю границу космического возраста}. В рамках YUCT наша наблюдаемая Вселенная — лишь один из многих локальных фазовых переходов внутри вечно вращающегося космоса, что естественно объясняет столь огромный период.

\subsubsection{Иерархия барионной массы и алгебраическая петля}

Одним из наиболее поразительных результатов Приложения $\Omega$ является установление прямой количественной связи между полной барионной массой Вселенной и массами протона и электрона. Эта связь не содержит свободных параметров и опирается исключительно на универсальные константы алгебраической петли YUCT ($S_{\mathrm{odd}}=6/5$, $S_{\mathrm{even}}=4/5$) и планковский масштаб:
\begin{equation}
	\boxed{ \frac{M_{\mathrm{bar}}}{m_p} = \left( \frac{M_{\mathrm{Pl}}}{m_e} \right)^{\mathcal{S}} },
	\label{eq:baryon_hierarchy_combined}
\end{equation}
где показатель
\begin{equation}
	\mathcal{S} \equiv S_{\mathrm{odd}} + S_{\mathrm{even}} + \frac{S_{\mathrm{odd}}}{S_{\mathrm{even}}} = \frac{6}{5} + \frac{4}{5} + \frac{3}{2} = 3.5.
\end{equation}
Численно, $(M_{\mathrm{Pl}}/m_e)^{3.5} \approx 1.88 \times 10^{78}$, что в пределах $\sim 20\%$ совпадает с наблюдаемым отношением $M_{\mathrm{bar}}/m_p \approx (1.4-1.7) \times 10^{78}$ (в зависимости от значения $H_0$). Это совпадение охватывает 78 порядков величины и связывает квантовые, гравитационные и космологические масштабы в единую безразмерную иерархию.

\subsubsection{Критическая масса локального Большого взрыва}

YUCT предсказывает, что наша Вселенная возникла как локальный координационный фазовый переход — «Большой взрыв» — внутри более обширного вращающегося космоса. Критическая масса предколлапсной области, инициирующей такой переход, строго фиксируется фундаментальными константами:
\begin{equation}
	M_{\mathrm{crit}} = 3 M_{\mathrm{Pl}} \approx 6.5 \times 10^{-8}\ \text{кг}.
	\label{eq:Mcrit_micro_combined}
\end{equation}
Это микроскопическое значение соответствует планковскому зародышу новой вселенной. Однако существует и макроскопический аналог — масса, при которой достаточно компактный астрофизический объект может претерпеть катастрофический фазовый переход, породив новую инфлирующую область внутри нашего пространства-времени. Этот макроскопический предел оказывается порядка полной массы наблюдаемой Вселенной:
\begin{equation}
	M_{\mathrm{crit}}^{\mathrm{(local)}} \approx M_{\mathrm{univ}} \approx 3 \times 10^{54}\ \text{кг}.
	\label{eq:Mcrit_macro_combined}
\end{equation}
Такое совпадение (с точностью до множителя порядка единицы) не случайно: оно является следствием того же универсального масштабного закона $K_{\mathrm{eff}} \propto R^{\beta d_f/2}$ с $\beta=2/3$ и фрактальной размерностью $d_f=2.2$, который управляет и рождением, и современной динамикой Вселенной.

\subsubsection{Связь с эмерджентной массой фотона}

Рассмотренное выше глобальное вращение и координационная природа тёмной энергии напрямую связаны с концепцией эмерджентной массы фотона (подраздел~\ref{subsec:time_photon_combined}). Вращение создаёт крупномасштабную завихренность, которая модулирует локальную координационную энтропию (Пралаю). Фотон, проходя через области с повышенной завихренностью, испытывает дополнительное увеличение $K_{\mathrm{eff}}$, что ведёт к временной эффективной массе. Таким образом, гравитационное линзирование, красное смещение и само ускоренное расширение Вселенной оказываются разными проявлениями единого координационного механизма, не требующего ни тёмной материи, ни тёмной энергии в их стандартном понимании.

% -------------------------------------------------------------------
\subsection{Предварительный характер и направления дальнейших исследований}

Концепции, представленные в этом разделе, были одними из центральных мотиваций на этапе формирования YUCT (2024–2025 гг.). Однако по мере получения большого массива количественных предсказаний (полуцелое квантование фермионных масс, периодическая таблица на основе $K_{\mathrm{eff}}$, законы фазовых переходов, биологические лагранжианы и т.д.) строгая математическая формулировка данных идей была отложена. Они остаются предварительными соображениями, подчёркивающими философскую целостность объединённой теории, и требуют дальнейшей проработки в качестве самостоятельных проверяемых предсказаний.

Будущие наблюдения (Euclid, SKA, CMB-S4, LISA) должны либо подтвердить предсказанные значения $\Omega_{\mathrm{rot}}$, $T_{\mathrm{rot}}$ и иерархию барионной массы, либо фальсифицировать теорию. Аналогично, эксперименты по гравитационному линзированию с высоким временным разрешением и лабораторные тесты модифицированного соотношения $E=mc^2$ способны проверить эффект эмерджентной массы фотона.

\vspace{0.5cm}
\noindent\textbf{Примечание:} Настоящий раздел включён только в русскоязычную версию документа и на данном этапе не входит в основные количественные выводы работы.

	\section{Заключение}
	
	Мы показали, что массы всех девяти заряженных фермионов Стандартной модели квантуются в полуцелых единицах $\ln q$, где $q = (3/2)^{1/3}$ — фундаментальный квант масштабирования данного подхода. Это описание использует всего 10 параметров (против 18 в традиционной параметризации) и достигает субпроцентной точности. Полуцелый шаг возникает из взаимодействия трёхмерного пространства и спиновой статистики. Статистическая значимость закономерности ($p < 10^{-15}$) сильно поддерживает координационное происхождение ароматов. Правило квантования даёт проверяемые предсказания для масс нейтрино и запрещает четвёртое поколение. Данная работа обеспечивает прочную эмпирическую основу для будущих первых принципов вывода масс фермионов в рамках теории.
	
	\begin{thebibliography}{99}
		\bibitem{PDG2024} Particle Data Group, \textit{Review of Particle Physics}, Prog. Theor. Exp. Phys. \textbf{2024}, 083C01 (2024).
		\bibitem{YUCT} A. V. Yakushev, \textit{Yakushev Unified Coordination Theory (YUCT) V2.0}, Zenodo, DOI:10.5281/zenodo.18444598 (2026).
	\end{thebibliography}
	
\end{document}